% MR57--0B
% Schock
% 14.0
% 2013-11-30

\label{m:am-schock}

\minisec{Allgemein}
\begin{enumerate}
  \item \TxMassVitMKBes
  \begin{compactitem}
    \item Situationsgerechte  Lagerung, je nach Diagnose und Schockart
    \item Monitoring: Besonders auf Haut, RR und HF achten!
  \end{compactitem}
  \item Beengende Kleidung öffnen
  \item Wärmeerhalt besonders beachten!
%   \item Spezielle Maßnahmen je nach Schockart bzw. Diagnose
  \item Transportentscheidung: Vorankündigung
  (Aviso); Spezialbett, je nach Schockart/Ursache
\end{enumerate}

\minisec{Speziell}

{\TableBaseModification
\begin{tabularx}{\linewidth}{XXXX}\addlinespace\toprule\addlinespace
\noindent\TabHeading{Schockart}
&
\noindent\TabHeading{Lagerung}
&
\noindent\TabHeading{Sonstige Maßnahmen}
&
\noindent\TabHeading{Transport}
\\\addlinespace\midrule\addlinespace
\noindent\TabSub{Volumenmangel}
&
Flachlagerung
&
Blutstillung

Schonende/schnelle Bergung\A{Schonend
  oder schnell?}
&
Schockraum, 

zügiger Abtransport
\\\addlinespace\midrule\addlinespace
\noindent\TabSub{Kardiogen}
&
Oberkörper hoch, evtl. Beine hängen  lassen.
&
Medikamentöse Therapie durch Notarzt\ZFootnote{(Diuretika, Katecholamine, Heparin)}
&
Ursachenabhängig
\\\addlinespace\midrule\addlinespace
\noindent\TabSub{Anaphylaktisch}
&
Flachlagerung, bei
Atemproblemen ggfs. stark erhöhten Oberkörper
&
Auslösendes Antigen nach Möglichkeit entfernen
&
Intensivstation
\\\addlinespace\midrule\addlinespace
\noindent\TabSub{Septisch, toxisch}
&
situationsgerecht, Beine hoch
&
&
Intensivstation
\\\addlinespace\midrule\addlinespace
\noindent\TabSub{Neurogen}
&
Flachlagerung, ggfs Immobilisation
(WS-Verletzung!)
&
Ursachenabhängig
\\\addlinespace\bottomrule
\end{tabularx}
}


\Customized{
\begin{description}
  \item[\CustAsboe] Beim anaphylaktischen Schock ist die Gabe
  von Adrenalin i.\,m. für NKV+ gem. Algorithmus 
  in der Arzneimittelliste 2 vorgesehen.
  \cite{Asb921Memo-AL-2011-000001}.
\end{description}
}
